\documentclass[notumble,combine]{leaflet}
%\documentclass[a4paper,notumble,combine,debug]{leaflet}
%\usepackage[utf8]{inputenc}
%\usepackage[T1]{fontenc}
%\usepackage{lmodern}
%\renewcommand{\bfdefault}{b}
%\usepackage{microtype}
\usepackage[dvipsnames,usenames]{color}
\usepackage{graphicx}
\usepackage{alltt}
\usepackage{url}
\usepackage{ascmac}
\usepackage{comment}
\usepackage{here}
\usepackage{wrapfig}
\usepackage{floatflt}
%\usepackage{ifthen}

\makeatletter
% Section color
%% from leaflet.cls
\renewcommand\section{\@startsection{section}{1}{\z@}%
  {-3.5ex \@plus -.75ex}%
  {1ex} %{1.5ex}%
  {\normalfont\large\sectfont\color{NavyBlue}}}
\renewcommand\subsection{\@startsection{subsection}{2}{\z@}%
  {-2.5ex plus -.5ex}%
  {1\p@} %{1ex}%
  {\normalfont\normalsize\sectfont\color{Green}}}
\makeatother

\graphicspath{{figures/}} 

\title{
	\includegraphics[width=\textwidth]{kbug-logo}\\
        \parbox[c]{2cm}{\resizebox{2cm}{!}{\textcolor{red}{\bf{@}}}}\\
	\includegraphics[width=\textwidth]{osc2015-kyoto-logo}\\
%	[\baselineskip]
	\vfill
	\parbox[c]{6cm}{\includegraphics[width=6cm]{kbug-summervacation-raspi}}\\
%	\parbox[c]{\textwidth}{\includegraphics[width=\textwidth]{kbug-summervacation-raspi}}\\
        \vfill
	\resizebox{7cm}{!}{\bf 関西*BSDユーザ会}\\
        \resizebox{6cm}{!}{{\bf \url{http://www.kbug.gr.jp/}}}\\
        \includegraphics[width=2cm]{kbug.gr.jp.eps}\\
        \resizebox{7cm}{!}{2015年8月7日(金),8日(土)}\\
}

\date{}
%\CutLine*{1} \CutLine*{2} \CutLine*{3} \CutLine*{4} \CutLine*{5}
%\CutLine*{6}

% \AddToBackground*{2}{% Background of a large page
%   \put(\LenToUnit{.25\paperwidth},\LenToUnit{.4\paperheight}){%
%     \includegraphics[width=14.85cm]{kbug-logo}
% }}


\begin{document}
\maketitle

\thispagestyle{empty}
\pagebreak{}
\section{関西*BSDユーザ会(K*BUG)ってなぁに?}
関西 *BSD ユーザ会 (Kansai *BSD Users Group; K*BUG)とは、
BSD系OSのユーザ同士の情報交換のための{\Large\em \textcolor{red}{場}}で、1999年に設立されました。

年に数度の勉強会やイベントの参加、そして{\Large\em \textcolor{red}{飲み会}}を行っています。

\subsection{K*BSDの基本理念}
K*BUGの基本理念はのとおりです。

\fbox{\begin{minipage}{\textwidth}
\begin{itemize}
\item 場の提供を目的とする
\item 人のケツは叩くが足は引っ張らない
\item 来るものは拒まず、猿ものは追わず
\item だれでも役員になれる。誰でも役員は止められる
\end{itemize}
\end{minipage}}

少し難しく感じるかもしれないですが、
\textcolor{red}{\bf\Large BSDへの愛と情熱}
があれば、
あなたが
\textcolor{blue}{やりたいと思うことができる場}
がK*BUGなのです。

%また、メンバーは色々な技術に詳しいですし、技術に関する話も大好きなので、
%あなたの疑問などもぶつけてみてください。

\section{K*BUG Keywords}
K*BUGでは、以下のようなKeywordに関する発表や展示を行いました。

{\large\bf\textcolor{red}{{BSD}}} (FreeBSD (PC-BSD), NetBSD, OpenBSD, DragonlyBSD, {\footnotesize Mac OS X(?), iOS(??) }, ...
の更新情報 / 利用方法 / ...,
arch {\footnotesize (i386, amd64, macppc, landisk, zaurus, wzero3, hpcmips, fonera, netwalker, ipaq, fonera, ...)},
kernel hackや新技術 (DTrace, ZFS, シリアルドライバ, ...),
FreeBSD portsってなに?,
Security, 
教育,
Prolog, Lisp, awk,
Squeak, Scratch,
OpenCV,
ContaoCMS, 
hardware (UPS, HDD, Server, Bluetooth, ...),
ものづくり,
電子工作,
ニコ動技術部 (猫耳サーボ % (USB audio servo controler)
, 鼻ホタル, ...),
雑誌付録基板 ,
Physical Computing (Gainer (gainerm-lib), Arduino, ...),
木彫, でーもんくん, でもんむし君 , \\
\hfill{\textcolor{red}{{\bf\large 飲み会}}}, and {\em\textcolor{red}{{ more with }}{\Large\bf\textcolor{red}{{YOU!!}}}}

\newpage
\section{K*BUGの活動}
最近の活動はこんな感じです。
もし、興味のあるテーマがあったら、
\begin{center}
{\huge\bf \textcolor{red}{遊びに来てね!!}}
\end{center}

\subsection{2015年6月27日(土)@京都 第3回研究会}
\begin{itemize}
\item Bluetooth LE
\item ゆるいpkgsrcの
\item LT的話
  \begin{itemize}
  \item linuxのBluetooth LEで相談
  \item raidframeと壊れたディスク、bootは不可問題
  \item Gdev: Open Source のGPGPU Runtime and Driver Software
  \item Dellの4KディスプレイがFreeBSDで使えへん問題
  \end{itemize}
\item 謎言語
\end{itemize}

\subsection{2015年5月16日(土)@京都 第2回研究会}
\begin{itemize}
\item 今さら sh
\end{itemize}

\subsection{2015年3月7日(土)@大阪 第1回研究会}
\begin{itemize}
\item うめきた9Fオフィスの紹介
\item blink1について
\item FreebSDでRaspberry  PIでGPIO
\item 周回遅れの HTTP/2
\item iPadに関する軽い話
\end{itemize}

\subsection{2015年1月24日(土)@京都\\\hfill 第16回定期総会 + 併設研究会}
\begin{itemize}
\item 高専におけるBSD化の試みとその顛末
\item dump(8)でバックアップ
\end{itemize}

\subsection{2014年8月23日(土)@京都 第3回研究会}
\begin{itemize}
\item wide cloud の FreeBSD 8.0 を freebsd-update(8) で upgrade
\item Bluetooth LE で遊びはじめた
\item tcpcrypt 導入してみた
\end{itemize}

\subsection{2014年6月28日(土)@大阪 第2回研究会}
\begin{itemize}
\item Wide Cloud で NetBSD VM が止まってしまって途方に暮れている話
\item 国土交通省CIMと OSS 
\item ソフトウェアラジオで遊ぼう (rtl2832)
\item ssh-agent を OS X の notification center に対応させる
\item HMZ-T2 + CubieBoard で OpenGL の wearable virtual reality experience
\end{itemize}

\subsection{2014年2月22日(土)@京都 第1回研究会}

\subsection{2013年12月7日(土)@大阪 \\\hfill 第15回定期総会+第6回研究会+忘年会}
\begin{itemize}
\item OSC Kansai@Kyoto 参加報告
\item AWK のあれこれ
\item Contao のイベント機能について
\item DKMSのはなし
\end{itemize}

\subsection{2013年10月19日(土)@大阪 第5回研究会}
\begin{itemize}
\item カーネル/VM 勉強会@関西 ごかいめ に合流
\end{itemize}

\subsection{2013年8月24日(土)@大阪 第4回研究会}
\begin{itemize}
\item *BSDでは生きにくい?
\item Wintel支配からの脱獄
\item rubygem and pkgsrc
\end{itemize}

%\newpage
\section{これからの関連イベント(予定)}
K*BUGでは、2ヶ月に1度程度の頻度で勉強会を行っています。
また、主に関西のイベントで、展示などを行っています。

現在、予定されているイベントは、以下のとおりです。
詳細は、K*BUGのWebページをご確認ください。

\begin{itemize} 
\item 2015年8月22日(土)@大阪 第4回研究会
\item 2015年10月24日(土)@京都 第5回研究会
\item 2015年11月6日(金),7日(土) KOF2015: 未定
\item 2015年12月5日(土)@京都\\
  \hfill 第17回定期総会 + 併設研究会
\end{itemize}

\section{写真集}
\subsection{KOF2012}
    \includegraphics[width=3.5cm]{KOF2012Booth}
    \includegraphics[width=4.5cm]{kof2012-mutoh}
\subsection{OSC2013 Kansai@Kyoto}
    \includegraphics[width=4cm]{osc2013kyoto}
    \includegraphics[width=4cm]{osc2013icat}

    \begin{center}
      \includegraphics[width=6cm]{MyKyotoIchiba}
    \end{center}
\newpage
\section{K*BUG at OSC2015 Kansai@Kyoto}
K*BUGは、今回以下のような活動で参加しています。

\subsection{チラシとシールの配布} % by 日本NetBSDユーザーグループ}
JNUGさんのブースをお借りして、このチラシとK*BUGシールを配布しています。

\subsection{JNUGセミナー:NetBSDのご紹介}
NetBSDは、自由に利用可能で、BSDライセンスに基づき再配布可能なUNIX-likeオペレーティングシステムです。
NetBSDの概要および最近の動向をご紹介します。

\begin{itemize}
\item 日時: 2015/8/7(金) 15:15-16:00
\item 会場: 会議室B
\item 発表: 蛯原 純 (The NetBSD Project)
\end{itemize}

\begin{center}
 \includegraphics[width=4cm]{oval-seal-4}
\end{center}

\vfill
\begin{minipage}{\textwidth}
\begin{boxnote}

\section{kbug-usersメーリングリスト}
K*BUGでは、K*BUGメンバーの情報交換や、イベントなどの情報伝達用に
kbug-usersメーリングリストを用意しています。

K*BUGメンバーは、基本的にはこのメーリングリストを読んでいることが期待
されます。

購読は、% 右のQRコードを使うか、
このURL
\footnote{\scriptsize \url{http://www.kbug.gr.jp/kbug-mailing-lists.html}}
をご参照ください。

\end{boxnote}
\end{minipage}
\vfill
\hfill
{\scriptsize \verb+$Date: 2015/08/06 04:38:37 $ $Revision: 1.29 $+}  
\end{document}
